\documentclass[10pt]{article}

\usepackage[margin=0.72in]{geometry}
\usepackage{booktabs}
\usepackage{caption}
\usepackage{enumitem}
\usepackage{float}
\usepackage[T1]{fontenc}
\usepackage{graphicx}
\usepackage{hyperref}
\usepackage{microtype}
\usepackage{tabularx}
\usepackage{xcolor}

\hypersetup{
  colorlinks=true,
  linkcolor=blue!55!black,
  urlcolor=blue!55!black,
  citecolor=blue!55!black
}

\setlength{\parindent}{0pt}
\setlength{\parskip}{0.55em}
\setlength{\emergencystretch}{2em}
\setlist[itemize]{leftmargin=1.35em, itemsep=0.22em, topsep=0.25em}
\captionsetup{font=small, labelfont=bf}

\newcommand{\code}[1]{\texttt{#1}}
\newcommand{\splitkey}[1]{\texttt{#1}}

\title{\vspace{-1.2em}\textbf{WindsorML Dataset Splits}}
\author{}
\date{}

\begin{document}
\maketitle
\vspace{-2.0em}

Deterministic train/validation/test splits for the
\href{https://huggingface.co/datasets/neashton/windsorml}{WindsorML dataset}
\cite{windsorml_dataset}. WindsorML contains 355 Halton-sampled variants of the
Windsor body, indexed from \code{run\_0} through \code{run\_354}. The associated
paper reports a 60/20/20 partition for its preliminary ML evaluation but does
not publish an exact case-membership list \cite{windsorml_paper}. For consistency
with the AhmedML and DrivAerML split packages, this package instead defines an
approximately 80/10/10 seed-42 baseline and companion data-efficiency and
out-of-distribution (OOD) splits.

The source of truth is \code{splits/manifest.json}, a flat JSON object with keys
such as \code{full\_train}, \code{full\_val}, and \code{geometry\_test}. This
repository intentionally stores only the manifest, small CSV inputs and metrics,
scripts, figures, and this report. Large STL and PNG assets belong in a sibling
asset directory.

\section*{Splits at a glance}

\begin{tabularx}{\textwidth}{@{}l l r r r X@{}}
\toprule
Split & Type & Train & Val & Test & What it tests \\
\midrule
\splitkey{full} & In-dist & 284 & 35 & 36 & Seed-42 random baseline, approximately 80/10/10 \\
\splitkey{medium} & In-dist & 95 & 35 & 36 & Data efficiency, 1/3 of \splitkey{full} training data \\
\splitkey{scarce} & In-dist & 47 & 35 & 36 & Data efficiency, 1/6 of \splitkey{full} training data \\
\splitkey{super\_scarce} & In-dist & 8 & 35 & 36 & Extreme data efficiency, 1/36 of \splitkey{full} training data \\
\splitkey{geometry} & OOD & 248 & 36 & 71 & STL-shape extrapolation using the top 20\% local Chamfer-isolation score \\
\splitkey{high\_drag} & OOD & 248 & 36 & 71 & High-drag extrapolation using the top 20\% fixed-reference \code{cd} \\
\splitkey{low\_drag} & OOD & 248 & 36 & 71 & Low-drag extrapolation using the bottom 20\% fixed-reference \code{cd} \\
\splitkey{image\_wake} & OOD & 248 & 36 & 71 & Low-speed wake extrapolation from near-centreline and near-base velocity PNGs \\
\bottomrule
\end{tabularx}

\section*{Which split should I use?}

\begin{itemize}
  \item \textbf{Standard baseline}: \splitkey{full}
  \item \textbf{Data efficiency}: compare \splitkey{super\_scarce}, \splitkey{scarce}, \splitkey{medium}, and \splitkey{full}
  \item \textbf{Shape extrapolation}: \splitkey{geometry}
  \item \textbf{Aerodynamic-coefficient extrapolation}: \splitkey{high\_drag} or \splitkey{low\_drag}
  \item \textbf{Wake-regime extrapolation}: \splitkey{image\_wake}
\end{itemize}

\section*{Using the committed splits}

For the standard use case, read \code{splits/manifest.json}; no STL, PNG, or
regeneration step is required. Each value is a sorted list of dataset directory
names.

\begin{verbatim}
import json
from pathlib import Path

manifest = json.loads(Path("splits/manifest.json").read_text())

train_ids = manifest["full_train"]
val_ids = manifest["full_val"]
test_ids = manifest["full_test"]
\end{verbatim}

Change the \code{full} prefix to \code{medium}, \code{scarce},
\code{super\_scarce}, \code{geometry}, \code{high\_drag}, \code{low\_drag}, or
\code{image\_wake} to select another split.

\section*{Design principles}

\begin{itemize}
  \item \textbf{Validation remains in-distribution with training.} For each OOD split, the test set is the held-out extreme 20\%; validation is the companion-package 10\% sample drawn only from the complementary population.
  \item \textbf{Companion-package ratios.} \splitkey{full} is approximately 80/10/10. OOD families are approximately 70/10/20 so the extreme test regime remains large enough to evaluate separately.
  \item \textbf{Nested data-efficiency subsets.} The 8-case, 47-case, and 95-case training sets are strict nested subsets of \splitkey{full\_train}; all share \splitkey{full\_val} and \splitkey{full\_test}.
  \item \textbf{Use the appropriate source quantity.} Drag splits use the constant-reference-area coefficients in \code{force\_mom\_all.csv}; geometry uses STL-surface Chamfer distance; image wake uses published streamwise-velocity PNGs.
  \item \textbf{Test-set integrity.} Test cases should not be used for normalization fitting, hyperparameter tuning, model selection, or repeated visual inspection during development.
\end{itemize}

Figure~\ref{fig:diagnostics} shows the exact train, validation, and test
membership against each split-defining quantity. The first panel confirms that
\splitkey{full} is random with respect to drag; the high- and low-drag panels
show the expected coefficient tails; the lower panels show the geometry and
image-wake OOD regions.

\begin{figure}[H]
  \centering
  \includegraphics[width=0.99\textwidth]{split_diagnostics.png}
  \caption{WindsorML split diagnostics. Points are colored by train, validation, and test membership. The panels show the full seed-42 baseline on \code{Cd}; high- and low-drag holdouts on \code{Cd}; the STL-Chamfer geometry holdout; the image-wake holdout on its low-speed score; and the image-wake membership on \code{Cd}.}
  \label{fig:diagnostics}
\end{figure}

\section*{Split details}

\subsection*{\splitkey{full}}

To match the AhmedML and DrivAerML split convention, \splitkey{full} uses an
approximately 80/10/10 partition: 284 train, 35 validation, and 36 test cases.
To make the unpublished membership reproducible, this repository constructs
\code{torch.randperm(355)} with \code{torch.Generator().manual\_seed(42)}, assigns
the first 284 entries to train, the next 35 to validation, and the final 36 to
test, then sorts each stored ID list. The resulting IDs are committed directly
in the generator, so PyTorch is not a runtime dependency.

\subsection*{\splitkey{medium}, \splitkey{scarce}, and \splitkey{super\_scarce}}

These families keep \splitkey{full\_val} and \splitkey{full\_test} fixed. Their
training sets are nested, greedy max-min subsets of \splitkey{full\_train} in
standardized \code{cd}, \code{cl}, and published aggregate geometry-feature
space. The procedure selects a force/geometry-extreme case first, then repeatedly
adds the candidate whose nearest selected neighbor is farthest away.

The WindsorML paper defines seven CAD variables: front-to-rear length ratio,
back-fast length ratio, nose-to-windshield height ratio, fast-back height ratio,
side taper, clearance, and bottom taper angle \cite{windsorml_paper}. The current
public \code{geo\_parameters\_all.csv} contains six of those variables plus
\code{frontal\_area}; \code{ratio\_length\_front\_rear} appears in per-run geometry
CSVs but is absent from the aggregate table. The nested selection uses the
columns actually present in the aggregate file. The \splitkey{geometry} OOD
split does not depend on that table omission because it is computed from STLs.

\subsection*{\splitkey{geometry}}

Each available \code{run\_N/windsor\_N.stl} is sampled with 4096 deterministic,
area-weighted surface points. Point clouds remain in the common dataset
coordinate frame and are scaled by the global median STL bounding-box diagonal.
The symmetric Chamfer RMS distance is calculated for every pair, and each run's
OOD score is its mean distance to the 10 nearest neighboring STLs. The 71
highest scores form \splitkey{geometry\_test}; 36 validation cases are sampled
from the remaining population, leaving 248 training cases.

The current highest observed geometry scores include \splitkey{run\_252},
\splitkey{run\_161}, \splitkey{run\_349}, \splitkey{run\_72}, and
\splitkey{run\_57}. Figure~\ref{fig:geometry_examples} compares the lowest and
highest observed score examples and overlays the complete side-view geometry.

\begin{figure}[H]
  \centering
  \includegraphics[width=0.98\textwidth]{geometry_score_examples.png}
  \caption{Geometry examples for \splitkey{geometry}. The low-score case is \splitkey{run\_304} with Chamfer score 0.01147, \code{Cd}=0.2823, and \code{Cl}=-0.2293. The high-score case is \splitkey{run\_252} with score 0.02076, \code{Cd}=0.3408, and \code{Cl}=0.8068. The lower panel superimposes both complete-car side views with transparent silhouettes.}
  \label{fig:geometry_examples}
\end{figure}

\subsection*{\splitkey{high\_drag} and \splitkey{low\_drag}}

WindsorML publishes both constant-reference-area and case-dependent-reference-area
force tables. These splits use \code{force\_mom\_all.csv}, whose constant reference
area makes \code{cd} directly comparable across geometry variants.
\splitkey{high\_drag} holds out the 71 largest values; \splitkey{low\_drag} holds
out the 71 smallest. Figure~\ref{fig:diagnostics} shows both tails.

The highest-drag runs include \splitkey{run\_303}, \splitkey{run\_192},
\splitkey{run\_346}, \splitkey{run\_10}, and \splitkey{run\_221}. The lowest-drag
runs include \splitkey{run\_238}, \splitkey{run\_125}, \splitkey{run\_306},
\splitkey{run\_133}, and \splitkey{run\_134}.

\subsection*{\splitkey{image\_wake}}

The paper states that the published images span 10 constant-\(z\) planes from
\(z=-0.4\) to \(0.4\) m and 80 constant-\(x\) planes from \(x=-0.5\) to
\(1.0\) m \cite{windsorml_paper}. The image-wake score uses
\code{view1\_constz} indices 4 and 5, which bracket the centreline, and
\code{view2\_constx} indices 53, 55, and 57, immediately downstream of the
Windsor base at \(x=0.48\) m. In fixed wake crops, blue/purple pixels indicate
lower streamwise velocity than the orange freestream. The score averages the
low-speed area and color intensity across the centreline and near-base views.
The 71 largest scores form \splitkey{image\_wake\_test}.

The highest observed image-wake scores include \splitkey{run\_8},
\splitkey{run\_315}, \splitkey{run\_57}, \splitkey{run\_332}, and
\splitkey{run\_329}. Figure~\ref{fig:wake_examples} shows why the score separates
the selected examples: the high-score case has a substantially larger low-speed
region in both views.

\begin{figure}[H]
  \centering
  \includegraphics[width=0.99\textwidth]{wake_score_examples.png}
  \caption{Image-wake examples for \splitkey{image\_wake}. The low-score case is \splitkey{run\_161} with score 0.0076, \code{Cd}=0.3226, and \code{Cl}=0.8371. The high-score case is \splitkey{run\_8} with score 0.1519, \code{Cd}=0.3059, and \code{Cl}=-0.2052. The top row shows a near-centreline constant-\(z\) view; the bottom row shows the near-base \code{X-53} plane.}
  \label{fig:wake_examples}
\end{figure}

\section*{Current Hub asset coverage}

The aggregate CSVs contain all 355 runs. At the Hub revision used for this
package, per-run STLs and targeted PNGs were available for \code{run\_0} through
\code{run\_349}; the same assets for \code{run\_350} through \code{run\_354}
returned HTTP 404. Their five geometry and image scores are estimated from the
five nearest observed runs in standardized force/aggregate-geometry space.
Both \code{data/chamfer\_metrics.csv} and \code{data/image\_metrics.csv} include
an observed flag and the neighbor IDs for every estimate. Re-running the metric
scripts after those Hub files become available will replace the estimates with
direct observations.

\section*{Repeatability and transparency}

The committed manifest is intended for normal benchmark use. The commands below
are only for auditing or rebuilding the split definitions. Keep large files in
a sibling directory:

\begin{verbatim}
ASSET_ROOT=../windsorml_hf_assets

python3 scripts/download_hf_inputs.py --output-dir "$ASSET_ROOT" \
  --include-stls --include-wake-images --include-geometry-images \
  --workers 6 --allow-missing

python3 scripts/compute_chamfer_splits.py --data-root "$ASSET_ROOT" \
  --output-dir /tmp/windsorml_chamfer_4096 --samples 4096 \
  --workers 16 --sample-workers 2 --runs all --allow-missing
cp /tmp/windsorml_chamfer_4096/chamfer_metrics.csv \
  data/chamfer_metrics.csv

python3 scripts/compute_image_metrics.py --asset-root "$ASSET_ROOT" \
  --data-root data --output data/image_metrics.csv
python3 scripts/generate_splits.py
python3 scripts/visualize_splits.py
python3 scripts/create_example_figures.py --asset-root "$ASSET_ROOT"
latexmk -pdf -cd docs/README.tex
\end{verbatim}

The Chamfer script requires \code{numpy}, \code{scipy}, and \code{trimesh}; the
image and figure scripts require \code{numpy}, \code{Pillow}, and
\code{matplotlib}. The repository's \code{data/} directory should remain limited
to small CSV files.

The committed source artifacts are:

\begin{verbatim}
data/force_mom_all.csv
data/geo_parameters_all.csv
data/chamfer_metrics.csv
data/image_metrics.csv
<asset-root>/run_*/windsor_*.stl
<asset-root>/run_*/images/windsor_*.png
<asset-root>/run_*/images/velocityxavg/*.png
docs/split_diagnostics.png
docs/geometry_score_examples.png
docs/wake_score_examples.png
splits/manifest.json
\end{verbatim}

\section*{Manifest format}

\begin{verbatim}
{
    "full_train": ["run_0", "run_1", "..."],
    "full_val": ["run_8", "..."],
    "full_test": ["run_7", "..."],
    "geometry_train": ["run_0", "..."]
}
\end{verbatim}

Case IDs match the on-disk dataset directory names and are sorted numerically.

{\small
\begin{thebibliography}{9}
\bibitem{windsorml_dataset}
WindsorML dataset. \url{https://huggingface.co/datasets/neashton/windsorml}.

\bibitem{windsorml_paper}
Ashton, N. et al. ``WindsorML: High-Fidelity Computational Fluid Dynamics Dataset for Automotive Aerodynamics.'' \url{https://arxiv.org/abs/2407.19320}.
\end{thebibliography}
}

\end{document}
